\begin{remark*}[Exposition]
Open and closed sets describe how subsets sit inside a metric space.  Open
sets are subsets that contain a small ball around each of their points.
Closed sets are subsets whose complements are open.
\end{remark*}

\begin{definitionbox}{Definition --- Metric Interior Point}
\begin{definition}[Metric Interior Point]\label{def:metric-interior-point}
Let \((X,d)\) be a metric space, let \(E\subseteq X\), and let \(x\in X\).
The point \(x\) is an \emph{interior point of \(E\)} if there exists
\(\delta>0\) such that
\[
  B_d(x;\delta)\subseteq E.
\]
\end{definition}
\end{definitionbox}

\begin{remark*}[Standard quantified statement]
\[
x\text{ is an interior point of }E
\quad\Longleftrightarrow\quad
x\in X
\;\land\;
\exists\delta>0{:}\ B_d(x;\delta)\subseteq E.
\]
\end{remark*}

\begin{remark*}[Interpretation]
An interior point of \(E\) is a point with positive-radius room inside \(E\).
The entire ball around the point must remain in \(E\), not merely the point
itself.
\end{remark*}

\begin{dependencies}
\begin{itemize}
  \item \hyperref[def:metric-space]{Definition: Metric Space}
  \item \hyperref[def:metric-open-ball]{Definition: Open Ball}
\end{itemize}
\end{dependencies}

\begin{definitionbox}{Definition --- Metric Interior}
\begin{definition}[Metric Interior]\label{def:metric-interior}
Let \((X,d)\) be a metric space and let \(E\subseteq X\).  The \emph{interior
of \(E\)} is the set of all interior points of \(E\).  It is denoted by
\[
  E^\circ.
\]
\end{definition}
\end{definitionbox}

\begin{remark*}[Standard quantified statement]
\[
E^\circ
=
\{x\in X:\exists\delta>0\text{ such that }B_d(x;\delta)\subseteq E\}.
\]
\end{remark*}

\begin{remark*}[Interpretation]
The interior \(E^\circ\) collects exactly the points of \(E\) that have
positive-radius room inside \(E\).
\end{remark*}

\begin{remark*}[Notation]
The notation \(E^\circ\) is read as the \emph{interior of \(E\)} or
\emph{\(E\)-interior}.
\end{remark*}

\begin{dependencies}
\begin{itemize}
  \item \hyperref[def:metric-space]{Definition: Metric Space}
  \item \hyperref[def:metric-interior-point]{Definition: Metric Interior Point}
\end{itemize}
\end{dependencies}

\begin{definitionbox}{Definition --- Metric Open Set}
\begin{definition}[Metric Open Set]\label{def:metric-open-set}
Let \((X,d)\) be a metric space.  A subset \(E\subseteq X\) is \emph{open in
\((X,d)\)} if either \(E=\varnothing\) or every \(x\in E\) is an interior point
of \(E\).
\end{definition}
\end{definitionbox}

\begin{remark*}[Standard quantified statement]
\[
E\text{ is open in }(X,d)
\quad\Longleftrightarrow\quad
E\subseteq X
\;\land\;
\bigl(E=\varnothing
\;\lor\;
\forall x\in E\,\exists\varepsilon>0{:}\ B_d(x;\varepsilon)\subseteq E\bigr).
\]
\end{remark*}

\begin{remark*}[Interpretation]
An open set contains room around each of its points.  The amount of room may
depend on the point, but every point of the set must have some positive-radius
ball that stays inside the set.
\end{remark*}

\begin{dependencies}
\begin{itemize}
  \item \hyperref[def:metric-space]{Definition: Metric Space}
  \item \hyperref[def:metric-open-ball]{Definition: Open Ball}
  \item \hyperref[def:metric-epsilon-neighborhood]{Definition: Epsilon Neighborhood}
  \item \hyperref[def:metric-interior-point]{Definition: Metric Interior Point}
\end{itemize}
\end{dependencies}

\begin{theorembox}{Theorem}
\begin{theorem}[Metric Open Set Closure Properties]
\label{thm:metric-open-set-closure-properties}
Let \(X\) be a metric space.  Then:
\begin{enumerate}
  \item the empty set \(\varnothing\) and the space \(X\) are open sets;
  \item an arbitrary union of open subsets of \(X\) is open;
  \item any finite intersection of open subsets of \(X\) is open.
\end{enumerate}

\smallskip
\noindent
\hyperref[prf:metric-open-set-closure-properties]{\textit{Go to proof.}}
\end{theorem}
\end{theorembox}

\begin{remark*}[Standard quantified statement]
If \(X\) is a metric space, then
\[
\varnothing\text{ is open in }X,
\qquad
X\text{ is open in }X,
\]
\[
\{U_\alpha\}_{\alpha\in A}\text{ open in }X
\Longrightarrow
\bigcup_{\alpha\in A}U_\alpha\text{ is open in }X,
\]
and, for \(n\in\mathbb{N}\),
\[
U_1,\ldots,U_n\text{ open in }X
\Longrightarrow
\bigcap_{k=1}^{n}U_k\text{ is open in }X.
\]
\end{remark*}

\begin{remark*}[Interpretation]
Metric open sets are stable under the same basic operations that define a
topology: arbitrary unions and finite intersections.  Thus every metric
determines a topology by declaring its metric-open subsets to be open.
\end{remark*}

\begin{remark*}[Historical note]
This result corresponds to Theorem~3.2 in Kainth's
\emph{A Comprehensive Textbook on Metric Spaces}
\cite{KainthComprehensiveMetricSpaces2023}.
\end{remark*}

\begin{dependencies}
\begin{itemize}
  \item \hyperref[def:metric-space]{Definition: Metric Space}
  \item \hyperref[def:metric-open-set]{Definition: Metric Open Set}
  \item \hyperref[def:metric-interior-point]{Definition: Metric Interior Point}
\end{itemize}
\end{dependencies}

\begin{theorembox}{Theorem}
\begin{theorem}[Open Balls are Open]
\label{thm:metric-open-balls-are-open}
Open balls in metric spaces are open sets.

\smallskip
\noindent
\hyperref[prf:metric-open-balls-are-open]{\textit{Go to proof.}}
\end{theorem}
\end{theorembox}

\begin{remark*}[Standard quantified statement]
\[
\forall (X,d)\,\forall x\in X\,\forall r>0{:}\quad
B_d(x;r)\text{ is open in }(X,d).
\]
\end{remark*}

\begin{remark*}[Interpretation]
Every metric ball contains room around each of its points.  The triangle
inequality is the mechanism that transfers room around the center to room
around an arbitrary point inside the ball.
\end{remark*}

\begin{remark*}[Historical note]
This result corresponds to Theorem~3.4 in Kainth's
\emph{A Comprehensive Textbook on Metric Spaces}
\cite{KainthComprehensiveMetricSpaces2023}.
\end{remark*}

\begin{dependencies}
\begin{itemize}
  \item \hyperref[def:metric-space]{Definition: Metric Space}
  \item \hyperref[def:metric-open-ball]{Definition: Open Ball}
  \item \hyperref[def:metric-open-set]{Definition: Metric Open Set}
\end{itemize}
\end{dependencies}

\begin{theorembox}{Theorem}
\begin{theorem}[Interior is the Largest Open Subset]
\label{thm:metric-interior-largest-open-subset}
Let \(E\) be a subset of a metric space \(X\).  Then \(E^\circ\) is the
largest open subset of \(X\) contained in \(E\).  In other words:
\begin{enumerate}
  \item \(E^\circ\subseteq E\);
  \item \(E^\circ\) is open;
  \item \(E^\circ\supseteq O\), for every open set \(O\subseteq E\).
\end{enumerate}

\smallskip
\noindent
\hyperref[prf:metric-interior-largest-open-subset]{\textit{Go to proof.}}
\end{theorem}
\end{theorembox}

\begin{remark*}[Standard quantified statement]
\[
E^\circ\subseteq E,
\qquad
E^\circ\text{ is open in }X,
\qquad
\forall O\subseteq E{:}\;
O\text{ open in }X\Longrightarrow O\subseteq E^\circ.
\]
\end{remark*}

\begin{remark*}[Interpretation]
The interior operation extracts the open part of a set.  It keeps exactly the
points with room inside \(E\), and it contains every open subset already lying
inside \(E\).
\end{remark*}

\begin{remark*}[Historical note]
This result corresponds to Theorem~3.5 in Kainth's
\emph{A Comprehensive Textbook on Metric Spaces}
\cite{KainthComprehensiveMetricSpaces2023}.
\end{remark*}

\begin{dependencies}
\begin{itemize}
  \item \hyperref[def:metric-space]{Definition: Metric Space}
  \item \hyperref[def:metric-interior]{Definition: Metric Interior}
  \item \hyperref[def:metric-open-set]{Definition: Metric Open Set}
  \item \hyperref[thm:metric-open-balls-are-open]{Theorem: Open Balls are Open}
\end{itemize}
\end{dependencies}

\begin{definitionbox}{Definition --- Metric Closed Set}
\begin{definition}[Metric Closed Set]\label{def:metric-closed-set}
Let \((X,d)\) be a metric space.  A subset \(E\subseteq X\) is \emph{closed in
\((X,d)\)} if its complement \(E^c=X\setminus E\) is open in \((X,d)\).
\end{definition}
\end{definitionbox}

\begin{remark*}[Standard quantified statement]
\[
E\text{ is closed in }(X,d)
\quad\Longleftrightarrow\quad
E\subseteq X
\;\land\;
E^c=X\setminus E
\;\land\;
E^c\text{ is open in }(X,d).
\]
\end{remark*}

\begin{remark*}[Interpretation]
A closed set is defined by the openness of what lies outside it.  In metric
spaces this means every point outside \(E\) has some positive-radius ball that
misses \(E\).
\end{remark*}

\begin{remark*}[Notation]
When \(E\) is a subset of a fixed ambient space \(X\), the symbol \(E^c\)
denotes the complement of \(E\) in \(X\):
\[
  E^c=X\setminus E.
\]
\end{remark*}

\begin{dependencies}
\begin{itemize}
  \item \hyperref[def:metric-space]{Definition: Metric Space}
  \item \hyperref[def:metric-open-set]{Definition: Metric Open Set}
\end{itemize}
\end{dependencies}
